\section{Background}


\subsection{Web browser requests}
Upon visiting a web page, the browser will make various requests to fetch embedded resources such as scripts, style sheets and images.
Depending on the relation between the embedding website and the site that the resources are hosted on, these can be \textit{same-origin}, \textit{same-site} or \textit{cross-site}.
If the resource shares the same scheme (i.e.\ http or https), host (e.g.\ www.example.com) and port (e.g.\ 80 or 443) as the embedding site, it is considered same-origin.
In case there is no exact match for the host, but the resource is located on the same registrable domain name, the effective top level domain plus one (\textit{eTLD+1}), as the embedding website (e.g.\ www.example.com and foo.example.com), it is considered same-site.
Finally, resources that have a different eTLD+1 domain with regard to the including website are considered cross-site, i.e., resources from \textit{tracker.com} included on \textit{example.com} are cross-site.

Prior to making the connection to the server, the domain name first needs to be resolved to an IP address.
In the most straightforward case, the DNS resolution of the domain name returns an \texttt{A} record containing the IP address.
However, the domain could also use a \texttt{CNAME} record to refer to any other domain name.
This can be an iterative process as the new domain name can again resolve to another \texttt{CNAME} record; this process continues until an \texttt{A} record is found.
Through this indirection of CNAMEs, the host that the browser connects to may belong to a different party, such as a tracker, than the domain it actually requests the resource from.
This means that requests to \textit{xxx.example.com} may actually be routed to a different site, such as \textit{yyy.tracker.com}.

\textbf{Cookie scoping}
Before a request is sent, the browser will first determine which cookies to attach in the HTTP request. 
This includes all cookies that were set on the same (sub)domain as the one where the request will be sent to.
Other cookies that will be included are those that were set by a same-site resource, i.e.\ either on another subdomain, or on the top domain, and had the \texttt{Domain} attribute set to the top domain, for instance by the following response header on https://sub.example.com/
\begin{lstlisting}[basicstyle=\ttfamily\footnotesize,]
Set-Cookie: cookie=value; Domain=example.com
\end{lstlisting}
Cookies that were set without the \texttt{Domain} attribute will only be included on requests that are same-origin to the response containing the \texttt{Set-Cookie} header.
The \texttt{SameSite} attribute on cookies determines whether a cookie will be included if the request is cross-site.
If the value of this attribute is set to \texttt{None}, no restrictions will be imposed; if it is set to \texttt{Lax} or \texttt{Strict}, it will not be included on requests to resources that are cross-site to the embedding website; the latter imposes further restrictions on top-level navigational requests.
Several browser vendors intend to move to a configuration that assigns \texttt{SameSite=Lax} to all cookies by default~\cite{chromiumsamesite,draft-west-cookie-incrementalism-01,firefoxsamesitelaxdefault}.
As such, for third-party tracking to continue to work, the cookies set by the trackers explicitly need to set the \texttt{SameSite=None} attribute, which may make them easier to distinguish.
For CNAME-based tracking, where the tracking requests are same-site, the move to SameSite cookie by default has no effect.

\subsection{Tracking}
\subsubsection{Third-party tracking}
In a typical tracking scenario, websites include resources from a third-party tracker in a cross-site context.
As a result, when a user visits one of these web pages, a cookie originating from the third party is stored in the visitor's browser.
The next time a user visits a website on which the same tracker is embedded, the browser will include the cookie in the request to the tracker.
This scheme allows trackers to identify users across different websites to build detailed profiles of their browsing behavior.
Such tracking has triggered privacy concerns and has resulted in substantial research effort to understand the complexity of the tracking ecosystem \cite{mayer2012third,englehardt2015cookies} and its evolution \cite{lerner2016internet}.


\subsubsection{First-party tracking}
In first-party tracking the script and associated analytics requests are loaded from a same-site origin.
Consequently, any cookie that is set will only be included with requests to the same site.
Historically, one method that was used to bypass this limitation was cookie matching~\cite{olejnik2014selling}, where requests containing the current cookie are sent to a common third-party domain.
However, such scripts can be blocked by anti-tracking tools based on simple matching rules.
Instead, the technique covered in this work uses a delegation of the domain name, which circumvents the majority of anti-tracking mechanisms currently offered to users.

\subsubsection{CNAME-based tracking}

\textbf{General overview}
In the typical case of third-party tracking, a website will include a JavaScript file from the tracker, which will then report the analytics information by sending (cross-site) requests to the tracker domain.
With CNAME-based tracking, the same operations are performed, except that the domain that the scripts are included from and where the analytics data is sent to, is a subdomain of the website.
For example, the website \texttt{example.com} would include a tracking script from \texttt{track.example.com}, thus effectively appearing as same-site to the including website.
Typically, the subdomain has a CNAME record that points to a server of the tracker.
An overview of the CNAME-based tracking scheme is shown in Figure~\ref{fig:cname-scheme-overview}.

\begin{figure}
	\centering
	\includegraphics[width=\linewidth]{figures/cname-scheme-overview.pdf} 
	\caption{Overview of CNAME-based tracking.}
	\label{fig:cname-scheme-overview}
\end{figure}

\noindent \textbf{Bypassing anti-tracking measures}
The CNAME tracking scheme has direct implications for many anti-tracking mechanisms.
Because the requests to the tracking services are same-site (i.e.\ they point to the same eTLD+1 domain as the visited website), countermeasures that aim to block third-party cookies, such as Safari's ITP, are effectively circumvented.
Other popular anti-tracking mechanisms that rely on blocking requests or cookies by using block lists (such as EasyPrivacy \cite{easyprivacy} or Disconnect.me \cite{disconnectme}) become much harder to maintain  when trackers are served from a custom subdomain that is unique to every website. 
To block CNAME-based tracking, block lists would need to contain an entry for every website that uses the CNAME-based tracking service, instead of a single entry per tracker or match all DNS-level domains, leading to greater performance costs.

As a consequence of how the CNAME-based tracking scheme is constructed, it faces certain limitations in comparison to third-party tracking.
For instance, there no longer exists a common identifier shared across the different websites (in typical third-party tracking, the third-party cookie is responsible for this functionality).
Consequently, visits to different websites cannot be attributed to the same user using standard web development features.



